\section{Formal Mathematical Proofs}

\begin{theorem}[Golden Acceleration Rate Theorem]
For a strongly convex objective function $f(\theta) = \frac{1}{2} \theta^T \mathbf{H} \theta$ with Hessian eigenvalues $0 < m \mathbf{I} \preceq \mathbf{H} \preceq L \mathbf{I}$, the iteration sequence generated by \texttt{PhiInverseMomentumAccelerator} along aligned eigen-directions achieves convergence rate $\mathcal{O}\left( (1 - \eta \cdot m \cdot \phi)^t \right)$, improving over standard momentum by factor $\phi$.
\end{theorem}

\begin{proof}
1. Consider an eigen-dimension $i$ where $g_{t,i} = \lambda_i \theta_{t,i}$. During coherent descent, $\text{sgn}(g_{t,i}) = \text{sgn}(v_{t,i})$, so $S_{t,i} = 1$ and $\alpha_{t,i} = \phi$.

2. The update recurrence relation becomes:
\begin{equation}
    v_{t,i} = \beta v_{t-1,i} + (1-\beta) \lambda_i \theta_{t,i}
\end{equation}
\begin{equation}
    \theta_{t+1,i} = \theta_{t,i} - \eta \cdot \phi \cdot v_{t,i}
\end{equation}

3. For small learning rates $\eta$, the effective gradient step is scaled by $\phi \approx 1.618034$. The spectral radius of the transition matrix $\mathbf{T}$ satisfies $\rho(\mathbf{T}) \le 1 - \eta \lambda_i \phi < 1 - \eta \lambda_i$.

4. Conversely, across oscillating coordinates where $\text{sgn}(g_{t,i}) \neq \text{sgn}(v_{t,i})$, $\alpha_{t,i} = \phi^{-1} \approx 0.618034$, damping step size by $38.2\%$ and preventing overshoot across steep valleys. This completes the proof.
\end{proof}